\chapter{Einleitung}

\section{Einführung und Relevanz}
Digitale Assistenzsysteme haben sich längst im digitalen beruflichen Alltag etabliert. Mit dem Aufkommen von immer leistungsstärkeren KI-Modellen, werden immer mehr datenverarbeitende Prozesse durch neue KI basierte Methodiken erweitert oder ersetzt. Effizienz, Automatisierbarkeit und Intelligenz werden oft als Mehrwert gepriesen und somit gerechtfertigt. Dazu wird dieser Trend durch die aktuelle wirtschaftliche Dynamik erheblich getrieben. KI wird als zentrale Komponente von zukünftiger Innovation, Wertschöpfung und Profitschöpfung positioniert. Dennoch wächst auch die gesellschaftliche Kritik am unreflektierten Einsatz von und Umgang mit diesen Technologien. Speziell der hohe Energieverbrauch der mit KI-Modellen verbundene Infrastruktur stehen in der Kritik. Dabei stehen ökologische und soziale Implikationen im Vordergrund. Vor allem der Einsatz von KI Modellen für Problemstellungen für die wesentlich effizientere Alternativen bisher eingesetzt wurden. Eine Suchanfrage über ein LLM verbraucht rund 10 mal so viel wie eine herkömmliche Google-Suche \cite{Electricity2024Analysisandforecastto2026}.  Es stellt sich die Frage ob der Einsatz von KI immer wirklich die sinnvollste technische und gesellschaftliche Lösung darstellt. 

\section{Ziel und Forschungsfrage}
Diese Arbeit untersucht die ethische Vertretbarkeit des Einsatzes von KI in datenverarbeitenden Assistenzsystemen, vor allem dann, wenn einfache und energieeffizientere Alternativen existieren.  Daraus ergibt sich folgende Forschungsfrage:

\begin{quote}
    Inwiefern ist der Einsatz von KI in datenverarbeitenden Assistenzsystemen ethisch vertretbar, wenn energieeffizientere und technisch einfachere Alternativen existieren? 
\end{quote}

Ziel der Arbeit ist Kriterien für die verantwortliche und nachhaltige Gestaltung technischer Systeme herauszuarbeiten, und anhand dieser die ethische Vertretbarkeit des Praxisprojekts auszuwerten.

\section{Bezug zum Praxisprojekt}
Die Überlegungen werden durch Analyse eines konkreten Praxisprojekts ergänzt. Im Rahmen des Praxisprojekts wird ein KI-basiertes Assistenzsystem entwickelt. Dieses Projekt dient als Fallstudie um ethische sowie technische Entscheidungen exemplarisch zu analysieren. Dieses bietet einen Einblick in reale technische und wirtschaftliche Entscheidung und wie diese ethisch reflektiert werden können.

Wenn ein KI-Modell Informationen generiert, deren Mehrwert gegenüber einer einfacheren Lösung marginal ist, dann rechtfertigt das den ökologischen und sozialen Aufwand kaum.